% =====================================================================
%  ICD — Interface Control Document — Bus CAN
%  Groupe 69 : Conception d'une Électronique Modulaire (Bus CAN) — ENSEA 2A
%
%  FORMULAIRE PDF REMPLISSABLE — v2
%  Compiler avec pdflatex (PAS xelatex/lualatex : les champs ne seraient pas générés).
%  Le PDF produit se remplit dans Acrobat Reader, Firefox, Edge, Aperçu (macOS)…
%
%  Un exemplaire par projet partenaire : modifier \Projet et \NumGroupe ci-dessous,
%  compiler, puis diffuser le PDF.
% =====================================================================
\documentclass[10pt,a4paper]{article}

% ---------- Encodage et langue ----------
\usepackage[utf8]{inputenc}
\usepackage[T1]{fontenc}
\usepackage[french]{babel}
\usepackage{lmodern}

% ---------- Mise en page ----------
\usepackage[margin=1.8cm]{geometry}
\usepackage{fancyhdr}
\usepackage{lastpage}
\usepackage{parskip}
\usepackage{needspace}

% ---------- Tableaux ----------
\usepackage{booktabs}
\usepackage{longtable}
\usepackage{tabularx}
\usepackage{array}
\usepackage{colortbl}

% ---------- Couleurs, schémas ----------
\usepackage{xcolor}
\usepackage{tikz}
\usetikzlibrary{positioning,calc}

% ---------- Formulaire PDF ----------
\usepackage{hyperref}
\usepackage{amssymb}   % \square, \boxtimes dans les lignes d'exemple

% =====================================================================
%  À MODIFIER POUR CHAQUE PROJET
% =====================================================================
\newcommand{\Projet}{[Nom du projet partenaire]}
\newcommand{\NumGroupe}{[N\textsuperscript{o}]}
\newcommand{\Version}{0.1}
\newcommand{\DateVersion}{17 septembre 2026}
\newcommand{\Statut}{Brouillon — première séance}

% ---------- Couleurs ----------
\definecolor{canblue}{HTML}{1F3864}
\definecolor{canamber}{HTML}{FFB627}
\definecolor{rowgrey}{HTML}{F2F2F2}
\definecolor{exgrey}{HTML}{666666}

\hypersetup{colorlinks=true, linkcolor=canblue, urlcolor=canblue}

% ---------- Macros de champs de formulaire ----------
% charsize=0pt : taille de police automatique, le texte rétrécit pour tenir dans la case.
\newcommand{\champ}[1]{%
  \TextField[name=#1, width=\linewidth, height=12pt, charsize=0pt,
    bordercolor={0.72 0.72 0.72}, backgroundcolor={0.965 0.965 0.965},
    borderwidth=0.4]{}}

% Champ sur deux lignes, pour le texte libre.
\newcommand{\champML}[1]{%
  \TextField[name=#1, multiline=true, width=\linewidth, height=24pt, charsize=7.5pt,
    bordercolor={0.72 0.72 0.72}, backgroundcolor={0.965 0.965 0.965},
    borderwidth=0.4]{}}

% Grand champ de commentaire libre.
\newcommand{\champBig}[2]{%
  \TextField[name=#1, multiline=true, width=\linewidth, height=#2, charsize=8pt,
    bordercolor={0.72 0.72 0.72}, backgroundcolor={0.965 0.965 0.965},
    borderwidth=0.4]{}}

\newcommand{\coche}[1]{%
  \CheckBox[name=#1, width=9pt, height=9pt,
    bordercolor={0.35 0.35 0.35}, backgroundcolor={0.965 0.965 0.965},
    borderwidth=0.5]{}}

% Texte des lignes d'exemple
\newcommand{\ex}[1]{\textcolor{exgrey}{\textit{#1}}}

% ---------- Bascule portrait / paysage SANS rotation de page ----------
% (une page tournée par pdflscape fait apparaître le texte des champs à 90°)
\newcommand{\paysage}{%
  \clearpage
  \pdfpagewidth=297mm \pdfpageheight=210mm
  \newgeometry{layoutwidth=297mm, layoutheight=210mm, margin=1.8cm}%
  \fancyhfoffset{0pt}}
\newcommand{\portrait}{%
  \clearpage
  \pdfpagewidth=210mm \pdfpageheight=297mm
  \restoregeometry
  \fancyhfoffset{0pt}}

% ---------- En-tête / pied ----------
\pagestyle{fancy}
\fancyhf{}
\lhead{\small ICD Bus CAN — \Projet}
\rhead{\small v\Version{} — \DateVersion}
\cfoot{\small Page \thepage\ / \pageref{LastPage}}
\renewcommand{\headrulewidth}{0.4pt}

\newcolumntype{L}[1]{>{\raggedright\arraybackslash}p{#1}}
\newcolumntype{C}[1]{>{\centering\arraybackslash}p{#1}}

% Un titre de section ne doit jamais être séparé de son contenu
\let\oldsection\section
\renewcommand{\section}{\Needspace*{10\baselineskip}\oldsection}
\let\oldsubsection\subsection
\renewcommand{\subsection}{\Needspace*{8\baselineskip}\oldsubsection}

% =====================================================================
\begin{document}
\begin{Form}


\begin{titlepage}
  \centering
  \vspace*{2cm}
  {\color{canblue}\Huge\bfseries Interface Control Document\par}
  \vspace{0.4cm}
  {\Large Bus CAN — \Projet\par}
  \vspace{0.3cm}
  {\color{canamber}\rule{0.5\textwidth}{2pt}\par}
  \vspace{2cm}

  \begin{tabular}{@{}r l@{}}
    \textbf{Version} & \Version \\[3pt]
    \textbf{Date} & \DateVersion \\[3pt]
    \textbf{Statut} & \Statut \\[3pt]
    \textbf{Projet partenaire} & \Projet{} (groupe \NumGroupe) \\[3pt]
    \textbf{Rédigé par} & Groupe 69 — Électronique Modulaire (Bus CAN) \\[3pt]
    & Samuel Gomri · Dylan Kiati · Hugo Mathoux \\[3pt]
    \textbf{Contact} & hugo.mathoux@ensea.fr \\[3pt]
    \textbf{Encadrants} & Yves Heckel · Julien Seigneurbieux \\
  \end{tabular}

  \vfill
  \begin{minipage}{0.88\textwidth}
    \small
    \textbf{Comment remplir ce document.} Les cases grisées sont des champs de saisie :
    cliquez dedans et tapez, puis \textbf{enregistrez} le PDF. Chaque section commence par
    un tableau d'exemples (en gris italique) montrant le niveau de détail attendu — il n'est
    pas modifiable, il sert uniquement de repère.

    \medskip
    Ce document est \textbf{vivant}. Il est normal qu'il soit incomplet en début de projet :
    remplissez ce que vous savez, laissez le reste vide. Chaque évolution donne lieu à une
    nouvelle version datée, tracée en section~\ref{sec:historique}.
  \end{minipage}

  \vspace{1.5cm}
  {\small ENSEA — Projets 2\textsuperscript{e} année — Module MANAGEMENT\_PROJET\_S7 — 2026--2027\par}
\end{titlepage}

\tableofcontents
\newpage


\section{Historique des versions et des rencontres}
\label{sec:historique}

Une ligne par rencontre entre le groupe 69 et l'équipe \Projet{}, et une ligne par version
diffusée. Si une rencontre produit une nouvelle version, la colonne Version l'indique ; sinon
elle reste vide. Ne jamais effacer une ligne : ce tableau sert de trace des décisions.

\subsection*{Exemple}
\begin{tabularx}{\textwidth}{L{2.2cm} C{1.4cm} L{3.4cm} X}
  \toprule
  \textbf{Date} & \textbf{Version} & \textbf{Participants} & \textbf{Décisions / modifications} \\
  \midrule
  \ex{29/09/2026} & \ex{0.2} & \ex{H. Mathoux, A. Dupont, B. Martin} & \ex{Ajout du capteur de température batterie (ID 0x301) ; débit confirmé à 500 kbit/s.} \\
  \ex{06/10/2026} & \ex{—} & \ex{H. Mathoux, A. Dupont} & \ex{Point d'avancement, pas de changement.} \\
  \bottomrule
\end{tabularx}

\subsection*{À remplir}
\begin{tabularx}{\textwidth}{L{2.2cm} C{1.4cm} L{3.4cm} X}
  \toprule
  \textbf{Date} & \textbf{Version} & \textbf{Participants} & \textbf{Décisions / modifications} \\
  \midrule
  \DateVersion & 0.1 & Groupe 69 + \Projet & Remise de la trame vierge, présentation du bus CAN. \\

  \champML{h1d} & \champML{h1v} & \champML{h1p} & \champML{h1m} \\
  \champML{h2d} & \champML{h2v} & \champML{h2p} & \champML{h2m} \\
  \champML{h3d} & \champML{h3v} & \champML{h3p} & \champML{h3m} \\
  \champML{h4d} & \champML{h4v} & \champML{h4p} & \champML{h4m} \\
  \champML{h5d} & \champML{h5v} & \champML{h5p} & \champML{h5m} \\
  \champML{h6d} & \champML{h6v} & \champML{h6p} & \champML{h6m} \\
  \champML{h7d} & \champML{h7v} & \champML{h7p} & \champML{h7m} \\
  \champML{h8d} & \champML{h8v} & \champML{h8p} & \champML{h8m} \\
  \champML{h9d} & \champML{h9v} & \champML{h9p} & \champML{h9m} \\
  \champML{h10d} & \champML{h10v} & \champML{h10p} & \champML{h10m} \\
  \champML{h11d} & \champML{h11v} & \champML{h11p} & \champML{h11m} \\
  \champML{h12d} & \champML{h12v} & \champML{h12p} & \champML{h12m} \\
  \bottomrule
\end{tabularx}


\section{Objet et périmètre}

\subsection{Objet du document}

Cet ICD définit l'interface CAN entre les sous-systèmes du projet \Projet{}. Pour chaque message
qui circule sur le bus, il précise : l'identifiant, l'émetteur, les récepteurs, le contenu et son
codage, la fréquence d'émission et la criticité. C'est la référence commune : une donnée non décrite
ici n'est pas censée circuler sur le bus.

\subsection{Ce que ce document n'est pas}

\begin{itemize}
  \item Ce n'est pas un contrat figé. Les besoins évoluent en cours de projet, c'est attendu.
        On met à jour, on incrémente la version, on prévient les personnes concernées.
  \item Ce n'est pas un cahier des charges technique du véhicule ou du robot : seul ce qui
        transite sur le bus CAN est décrit.
  \item Ce n'est pas un document du groupe 69 seul : il appartient aux deux équipes.
\end{itemize}

\subsection{Contacts de l'équipe \Projet}

Les membres de votre équipe et, si vous l'avez déjà décidé, qui s'occupe de quoi.
Le rôle peut rester vide.

\subsubsection*{Exemple}
\begin{tabularx}{\textwidth}{L{4.2cm} L{5cm} X}
  \toprule
  \textbf{Nom} & \textbf{Rôle dans le projet} & \textbf{Email} \\
  \midrule
  \ex{A. Dupont} & \ex{Électronique de puissance} & \ex{a.dupont@ensea.fr} \\
  \bottomrule
\end{tabularx}

\subsubsection*{À remplir}
\begin{tabularx}{\textwidth}{L{4.2cm} L{5cm} X}
  \toprule
  \textbf{Nom} & \textbf{Rôle dans le projet} & \textbf{Email} \\
  \midrule

  \champ{c1n} & \champ{c1r} & \champ{c1e} \\
  \champ{c2n} & \champ{c2r} & \champ{c2e} \\
  \champ{c3n} & \champ{c3r} & \champ{c3e} \\
  \champ{c4n} & \champ{c4r} & \champ{c4e} \\
  \champ{c5n} & \champ{c5r} & \champ{c5e} \\
  \champ{c6n} & \champ{c6r} & \champ{c6e} \\
  \bottomrule
\end{tabularx}


\section{Paramètres du bus}

Valeurs proposées par défaut par le groupe 69, identiques pour tous les projets. Cochez « OK » si
la proposition vous convient, sinon indiquez ce que vous préférez dans la dernière colonne.
En cas de doute, laissez notre proposition.

\begin{tabularx}{\textwidth}{L{4.0cm} L{6.6cm} C{1.0cm} X}
  \toprule
  \textbf{Paramètre} & \textbf{Notre proposition} & \textbf{OK} & \textbf{Si différent, préciser} \\
  \midrule

  \rowcolor{rowgrey} Version du protocole & CAN 2.0A (classique) & \coche{parok1} & \champ{parc1} \\
   Format des identifiants & 11 bits (standard) & \coche{parok2} & \champ{parc2} \\
  \rowcolor{rowgrey} Débit & 500 kbit/s & \coche{parok3} & \champ{parc3} \\
   Terminaisons & 120\,$\Omega$ aux deux extrémités du bus & \coche{parok4} & \champ{parc4} \\
  \rowcolor{rowgrey} Longueur max. de dérivation & 30 cm par nœud & \coche{parok5} & \champ{parc5} \\
   Câble & Paire torsadée 120\,$\Omega$ & \coche{parok6} & \champ{parc6} \\
  \rowcolor{rowgrey} Transceiver & Conforme ISO 11898-2 (interopérabilité garantie) & \coche{parok7} & \champ{parc7} \\
   Niveau logique TXD / RXD & \textbf{3,3 V} — la plupart des µC (STM32, ESP32) y sont & \coche{parok8} & \champ{parc8} \\
  \rowcolor{rowgrey} Alimentation du transceiver & 5 V + broche VIO en 3,3 V si le 5 V est dispo (MCP2562, TJA1051T/3) ; sinon transceiver 3,3 V (SN65HVD230) & \coche{parok9} & \champ{parc9} \\
  \bottomrule
\end{tabularx}

\medskip
\noindent\footnotesize\textbf{Pourquoi ces deux cas d'alimentation ?} Le bus CAN est interopérable
quelle que soit la tension d'alimentation du transceiver : un nœud 3,3 V et un nœud 5 V cohabitent
sans adaptateur sur le même bus, car la norme définit des seuils de détection très en dessous des
niveaux produits par les deux familles. Ce qui doit être commun, c'est le \textbf{niveau logique
TXD/RXD à 3,3 V}, pour ne pas envoyer 5 V dans une broche de microcontrôleur prévue pour 3,3 V.
Les anciens transceivers TJA1050 et MCP2551 ne conviennent pas : ils n'ont pas de broche VIO.
\normalsize

\medskip
\noindent\footnotesize\textit{Commentaire libre sur les paramètres du bus :}\\[2pt]
\champBig{parlibre}{40pt}
\normalsize


\section{Topologie et inventaire des nœuds}

\subsection{Schéma de principe}

Schéma générique : deux fils partagés, chaque nœud branché en dérivation, une terminaison
120\,$\Omega$ à chaque extrémité. À adapter au nombre réel de nœuds de votre projet.

\begin{center}
\begin{tikzpicture}[
    nd/.style={draw, rounded corners=2pt, fill=canblue, text=white, minimum width=2.1cm,
                 minimum height=0.75cm, font=\footnotesize\bfseries, align=center},
    term/.style={draw, fill=canamber, minimum width=0.22cm, minimum height=0.7cm, inner sep=0},
    bus/.style={line width=2pt},
    stub/.style={line width=1pt},
  ]
  \draw[bus, canamber] (0,0.3) -- (12,0.3);
  \draw[bus, canblue]  (0,0)   -- (12,0);
  \node[left, font=\footnotesize\bfseries] at (0,0.3) {CAN\_H};
  \node[left, font=\footnotesize\bfseries] at (0,0)   {CAN\_L};
  \node[term] at (0.12,0.15) {};
  \node[term] at (11.88,0.15) {};
  \node[above, font=\scriptsize] at (0.12,0.55) {120\,$\Omega$};
  \node[above, font=\scriptsize] at (11.88,0.55) {120\,$\Omega$};
  \foreach \x/\name in {1.6/{Nœud 1}, 4.2/{Nœud 2}, 6.8/{Nœud 3}, 9.4/{Nœud 4}} {
    \draw[stub, canblue] (\x,0) -- (\x,-0.6);
    \node[nd] at (\x,-1.0) {\name};
  }
\end{tikzpicture}
\end{center}

\subsection{Inventaire des nœuds}

Un nœud = tout équipement qui émet ou reçoit sur le bus. Un capteur qui ne parle pas CAN
nativement (analogique, I\textsuperscript{2}C, SPI…) passe par un microcontrôleur pont :
c'est ce microcontrôleur qui est le nœud. Un nœud peut à la fois émettre et recevoir :
cochez les deux cases.

\subsubsection*{Exemples}
\begin{tabularx}{\textwidth}{C{0.7cm} L{3.4cm} L{3.4cm} C{1.1cm} C{1.2cm} L{2.6cm} X}
  \toprule
  \textbf{N\textsuperscript{o}} & \textbf{Nœud} & \textbf{µC / transceiver} & \textbf{Émet} & \textbf{Reçoit} & \textbf{Responsable} & \textbf{Statut} \\
  \midrule
  \ex{1} & \ex{BMS batterie} & \ex{STM32F103 + MCP2562} & \ex{$\boxtimes$} & \ex{$\square$} & \ex{A. Dupont} & \ex{À faire} \\
  \ex{2} & \ex{Carte moteur} & \ex{STM32G431 + TJA1051T/3} & \ex{$\boxtimes$} & \ex{$\boxtimes$} & \ex{B. Martin} & \ex{En cours} \\
  \ex{3} & \ex{Tableau de bord} & \ex{ESP32 + SN65HVD230} & \ex{$\square$} & \ex{$\boxtimes$} & \ex{C. Durand} & \ex{Fait} \\
  \bottomrule
\end{tabularx}

\subsubsection*{À remplir}
\begin{tabularx}{\textwidth}{C{0.7cm} L{3.4cm} L{3.4cm} C{1.1cm} C{1.2cm} L{2.6cm} X}
  \toprule
  \textbf{N\textsuperscript{o}} & \textbf{Nœud} & \textbf{µC / transceiver} & \textbf{Émet} & \textbf{Reçoit} & \textbf{Responsable} & \textbf{Statut} \\
  \midrule

  1 & \champ{nd1n} & \champ{nd1u} & \coche{nd1e} & \coche{nd1r} & \champ{nd1p} & \champ{nd1s} \\
  2 & \champ{nd2n} & \champ{nd2u} & \coche{nd2e} & \coche{nd2r} & \champ{nd2p} & \champ{nd2s} \\
  3 & \champ{nd3n} & \champ{nd3u} & \coche{nd3e} & \coche{nd3r} & \champ{nd3p} & \champ{nd3s} \\
  4 & \champ{nd4n} & \champ{nd4u} & \coche{nd4e} & \coche{nd4r} & \champ{nd4p} & \champ{nd4s} \\
  5 & \champ{nd5n} & \champ{nd5u} & \coche{nd5e} & \coche{nd5r} & \champ{nd5p} & \champ{nd5s} \\
  6 & \champ{nd6n} & \champ{nd6u} & \coche{nd6e} & \coche{nd6r} & \champ{nd6p} & \champ{nd6s} \\
  7 & \champ{nd7n} & \champ{nd7u} & \coche{nd7e} & \coche{nd7r} & \champ{nd7p} & \champ{nd7s} \\
  8 & \champ{nd8n} & \champ{nd8u} & \coche{nd8e} & \coche{nd8r} & \champ{nd8p} & \champ{nd8s} \\
  9 & \champ{nd9n} & \champ{nd9u} & \coche{nd9e} & \coche{nd9r} & \champ{nd9p} & \champ{nd9s} \\
  10 & \champ{nd10n} & \champ{nd10u} & \coche{nd10e} & \coche{nd10r} & \champ{nd10p} & \champ{nd10s} \\
  11 & \champ{nd11n} & \champ{nd11u} & \coche{nd11e} & \coche{nd11r} & \champ{nd11p} & \champ{nd11s} \\
  12 & \champ{nd12n} & \champ{nd12u} & \coche{nd12e} & \coche{nd12r} & \champ{nd12p} & \champ{nd12s} \\
  \bottomrule
\end{tabularx}


\clearpage
\section{Plan d'allocation des identifiants}
\label{sec:ids}

Rien à remplir ici. Rappel : sur un bus CAN, \textbf{l'identifiant le plus faible a la priorité
la plus haute} lors de l'arbitrage. Ce découpage réserve les ID les plus bas aux messages critiques
et laisse de la marge dans chaque plage pour les ajouts futurs. Il sert à choisir les identifiants
de la section~\ref{sec:catalogue} — que le groupe 69 attribuera avec vous.

\begin{tabularx}{\textwidth}{C{2.6cm} L{3.8cm} X L{2.2cm}}
  \toprule
  \textbf{Plage (hex)} & \textbf{Catégorie} & \textbf{Exemples} & \textbf{Priorité} \\
  \midrule
  \rowcolor{rowgrey}
  0x000 -- 0x07F & Sécurité / urgence & Arrêt d'urgence, défaut batterie, perte de contrôle & Critique \\
  0x080 -- 0x0FF & Commandes temps réel & Consigne moteur, freinage & Très haute \\
  \rowcolor{rowgrey}
  0x100 -- 0x1FF & Mesures rapides & Vitesse, courant, IMU & Haute \\
  0x200 -- 0x2FF & Mesures lentes & Températures, tensions, niveaux & Moyenne \\
  \rowcolor{rowgrey}
  0x300 -- 0x3FF & États et diagnostics & État des nœuds, compteurs d'erreurs & Basse \\
  0x400 -- 0x4FF & IHM / affichage & Données pour le tableau de bord & Basse \\
  \rowcolor{rowgrey}
  0x500 -- 0x7FF & Réservé & Extensions futures, debug & --- \\
  \bottomrule
\end{tabularx}


\paysage
\setlength{\tabcolsep}{2.5pt}
\section{Catalogue des messages}
\label{sec:catalogue}

Une ligne par information qui doit circuler sur le bus. Commencez par « Information transportée »
et « Émetteur », puis Min / Max / Précision utile : ce sont ces trois colonnes qui permettent de
choisir le codage et le nombre d'octets. La colonne \textbf{ID} est remplie par le groupe 69, pour
éviter les collisions d'identifiants. Un message a \textbf{un seul émetteur}.

\subsection*{Exemples}
\begin{longtable}{C{1.3cm} L{2.9cm} L{2.3cm} L{2.3cm} L{3.4cm} L{1.7cm} C{1.3cm} C{1.3cm} L{1.9cm} C{1.5cm} L{1.7cm}}
  \toprule
  \textbf{ID} & \textbf{Nom du message} & \textbf{Émetteur} & \textbf{Récepteur(s)} & \textbf{Information transportée} & \textbf{Unité physique} & \textbf{Min} & \textbf{Max} & \textbf{Précision utile} & \textbf{Période} & \textbf{Criticité} \\
  \midrule
  \endhead
  \ex{0x010} & \ex{EMERG\_STOP} & \ex{BMS} & \ex{Tous} & \ex{Arrêt d'urgence} & \ex{---} & \ex{0} & \ex{1} & \ex{booléen} & \ex{100 ms} & \ex{Critique} \\
  \ex{0x100} & \ex{WHEEL\_SPEED} & \ex{Capteur vitesse} & \ex{IHM} & \ex{Vitesse roue} & \ex{km/h} & \ex{0} & \ex{80} & \ex{0,01 km/h} & \ex{50 ms} & \ex{Haute} \\
  \ex{0x200} & \ex{MOTOR\_TEMP} & \ex{Carte puissance} & \ex{IHM} & \ex{Température moteur} & \ex{$^\circ$C} & \ex{-20} & \ex{150} & \ex{0,1 $^\circ$C} & \ex{1 s} & \ex{Moyenne} \\
  \bottomrule
\end{longtable}

\subsection*{À remplir}
\begin{longtable}{C{1.3cm} L{2.9cm} L{2.3cm} L{2.3cm} L{3.4cm} L{1.7cm} C{1.3cm} C{1.3cm} L{1.9cm} C{1.5cm} L{1.7cm}}
  \toprule
  \textbf{ID} & \textbf{Nom du message} & \textbf{Émetteur} & \textbf{Récepteur(s)} & \textbf{Information transportée} & \textbf{Unité physique} & \textbf{Min} & \textbf{Max} & \textbf{Précision utile} & \textbf{Période} & \textbf{Criticité} \\
  \midrule
  \endhead

  \champ{m1id} & \champ{m1nom} & \champ{m1em} & \champ{m1re} & \champML{m1inf} & \champ{m1un} & \champ{m1min} & \champ{m1max} & \champ{m1pre} & \champ{m1per} & \champ{m1cri} \\
  \champ{m2id} & \champ{m2nom} & \champ{m2em} & \champ{m2re} & \champML{m2inf} & \champ{m2un} & \champ{m2min} & \champ{m2max} & \champ{m2pre} & \champ{m2per} & \champ{m2cri} \\
  \champ{m3id} & \champ{m3nom} & \champ{m3em} & \champ{m3re} & \champML{m3inf} & \champ{m3un} & \champ{m3min} & \champ{m3max} & \champ{m3pre} & \champ{m3per} & \champ{m3cri} \\
  \champ{m4id} & \champ{m4nom} & \champ{m4em} & \champ{m4re} & \champML{m4inf} & \champ{m4un} & \champ{m4min} & \champ{m4max} & \champ{m4pre} & \champ{m4per} & \champ{m4cri} \\
  \champ{m5id} & \champ{m5nom} & \champ{m5em} & \champ{m5re} & \champML{m5inf} & \champ{m5un} & \champ{m5min} & \champ{m5max} & \champ{m5pre} & \champ{m5per} & \champ{m5cri} \\
  \champ{m6id} & \champ{m6nom} & \champ{m6em} & \champ{m6re} & \champML{m6inf} & \champ{m6un} & \champ{m6min} & \champ{m6max} & \champ{m6pre} & \champ{m6per} & \champ{m6cri} \\
  \champ{m7id} & \champ{m7nom} & \champ{m7em} & \champ{m7re} & \champML{m7inf} & \champ{m7un} & \champ{m7min} & \champ{m7max} & \champ{m7pre} & \champ{m7per} & \champ{m7cri} \\
  \champ{m8id} & \champ{m8nom} & \champ{m8em} & \champ{m8re} & \champML{m8inf} & \champ{m8un} & \champ{m8min} & \champ{m8max} & \champ{m8pre} & \champ{m8per} & \champ{m8cri} \\
  \champ{m9id} & \champ{m9nom} & \champ{m9em} & \champ{m9re} & \champML{m9inf} & \champ{m9un} & \champ{m9min} & \champ{m9max} & \champ{m9pre} & \champ{m9per} & \champ{m9cri} \\
  \champ{m10id} & \champ{m10nom} & \champ{m10em} & \champ{m10re} & \champML{m10inf} & \champ{m10un} & \champ{m10min} & \champ{m10max} & \champ{m10pre} & \champ{m10per} & \champ{m10cri} \\
  \champ{m11id} & \champ{m11nom} & \champ{m11em} & \champ{m11re} & \champML{m11inf} & \champ{m11un} & \champ{m11min} & \champ{m11max} & \champ{m11pre} & \champ{m11per} & \champ{m11cri} \\
  \champ{m12id} & \champ{m12nom} & \champ{m12em} & \champ{m12re} & \champML{m12inf} & \champ{m12un} & \champ{m12min} & \champ{m12max} & \champ{m12pre} & \champ{m12per} & \champ{m12cri} \\
  \champ{m13id} & \champ{m13nom} & \champ{m13em} & \champ{m13re} & \champML{m13inf} & \champ{m13un} & \champ{m13min} & \champ{m13max} & \champ{m13pre} & \champ{m13per} & \champ{m13cri} \\
  \champ{m14id} & \champ{m14nom} & \champ{m14em} & \champ{m14re} & \champML{m14inf} & \champ{m14un} & \champ{m14min} & \champ{m14max} & \champ{m14pre} & \champ{m14per} & \champ{m14cri} \\
  \champ{m15id} & \champ{m15nom} & \champ{m15em} & \champ{m15re} & \champML{m15inf} & \champ{m15un} & \champ{m15min} & \champ{m15max} & \champ{m15pre} & \champ{m15per} & \champ{m15cri} \\
  \champ{m16id} & \champ{m16nom} & \champ{m16em} & \champ{m16re} & \champML{m16inf} & \champ{m16un} & \champ{m16min} & \champ{m16max} & \champ{m16pre} & \champ{m16per} & \champ{m16cri} \\
  \champ{m17id} & \champ{m17nom} & \champ{m17em} & \champ{m17re} & \champML{m17inf} & \champ{m17un} & \champ{m17min} & \champ{m17max} & \champ{m17pre} & \champ{m17per} & \champ{m17cri} \\
  \champ{m18id} & \champ{m18nom} & \champ{m18em} & \champ{m18re} & \champML{m18inf} & \champ{m18un} & \champ{m18min} & \champ{m18max} & \champ{m18pre} & \champ{m18per} & \champ{m18cri} \\
  \champ{m19id} & \champ{m19nom} & \champ{m19em} & \champ{m19re} & \champML{m19inf} & \champ{m19un} & \champ{m19min} & \champ{m19max} & \champ{m19pre} & \champ{m19per} & \champ{m19cri} \\
  \champ{m20id} & \champ{m20nom} & \champ{m20em} & \champ{m20re} & \champML{m20inf} & \champ{m20un} & \champ{m20min} & \champ{m20max} & \champ{m20pre} & \champ{m20per} & \champ{m20cri} \\
  \champ{m21id} & \champ{m21nom} & \champ{m21em} & \champ{m21re} & \champML{m21inf} & \champ{m21un} & \champ{m21min} & \champ{m21max} & \champ{m21pre} & \champ{m21per} & \champ{m21cri} \\
  \champ{m22id} & \champ{m22nom} & \champ{m22em} & \champ{m22re} & \champML{m22inf} & \champ{m22un} & \champ{m22min} & \champ{m22max} & \champ{m22pre} & \champ{m22per} & \champ{m22cri} \\
  \champ{m23id} & \champ{m23nom} & \champ{m23em} & \champ{m23re} & \champML{m23inf} & \champ{m23un} & \champ{m23min} & \champ{m23max} & \champ{m23pre} & \champ{m23per} & \champ{m23cri} \\
  \champ{m24id} & \champ{m24nom} & \champ{m24em} & \champ{m24re} & \champML{m24inf} & \champ{m24un} & \champ{m24min} & \champ{m24max} & \champ{m24pre} & \champ{m24per} & \champ{m24cri} \\
  \champ{m25id} & \champ{m25nom} & \champ{m25em} & \champ{m25re} & \champML{m25inf} & \champ{m25un} & \champ{m25min} & \champ{m25max} & \champ{m25pre} & \champ{m25per} & \champ{m25cri} \\
  \champ{m26id} & \champ{m26nom} & \champ{m26em} & \champ{m26re} & \champML{m26inf} & \champ{m26un} & \champ{m26min} & \champ{m26max} & \champ{m26pre} & \champ{m26per} & \champ{m26cri} \\
  \champ{m27id} & \champ{m27nom} & \champ{m27em} & \champ{m27re} & \champML{m27inf} & \champ{m27un} & \champ{m27min} & \champ{m27max} & \champ{m27pre} & \champ{m27per} & \champ{m27cri} \\
  \champ{m28id} & \champ{m28nom} & \champ{m28em} & \champ{m28re} & \champML{m28inf} & \champ{m28un} & \champ{m28min} & \champ{m28max} & \champ{m28pre} & \champ{m28per} & \champ{m28cri} \\
  \champ{m29id} & \champ{m29nom} & \champ{m29em} & \champ{m29re} & \champML{m29inf} & \champ{m29un} & \champ{m29min} & \champ{m29max} & \champ{m29pre} & \champ{m29per} & \champ{m29cri} \\
  \champ{m30id} & \champ{m30nom} & \champ{m30em} & \champ{m30re} & \champML{m30inf} & \champ{m30un} & \champ{m30min} & \champ{m30max} & \champ{m30pre} & \champ{m30per} & \champ{m30cri} \\
  \bottomrule
\end{longtable}

\noindent\footnotesize\textbf{Légende.} \textit{Période} : intervalle entre deux émissions, ou « événementiel » si le message n'est émis que sur changement.
\textit{Min / Max} : plage de valeurs physiques attendue. \textit{Précision utile} : plus petit écart qu'il faut pouvoir distinguer (ex. 0,1 $^\circ$C).
\textit{Criticité} : Critique / Haute / Moyenne / Basse — doit rester cohérente avec la plage d'ID (section~\ref{sec:ids}).
\normalsize


\clearpage
\section{Codage des messages}
\label{sec:codage}

C'est ici qu'on fixe le codage exact de chaque octet de chaque message : sans cette étape, deux
équipes peuvent être d'accord sur « vitesse en km/h » et pourtant coder l'une en dixièmes, l'autre
en centièmes. \textbf{À remplir dans un second temps, avec le groupe 69}, une fois le catalogue
(section~\ref{sec:catalogue}) stabilisé. Une ligne par signal ; un message qui transporte plusieurs
signaux occupe plusieurs lignes consécutives. Le nombre d'octets d'un message (DLC) se déduit de
la dernière ligne qui le concerne.

\subsection*{Conventions de codage}

Proposées par le groupe 69. Cochez « OK » ou précisez ce que vous préférez.

\begin{tabularx}{0.8\textwidth}{L{4.4cm} L{7.2cm} C{1.1cm} X}
  \toprule
  \textbf{Convention} & \textbf{Notre proposition} & \textbf{OK} & \textbf{Si différent, préciser} \\
  \midrule
  \rowcolor{rowgrey} Ordre des octets & Little-endian (octet de poids faible en premier) & \coche{cvok1} & \champ{cvc1} \\
  Entiers signés & Complément à deux & \coche{cvok2} & \champ{cvc2} \\
  \rowcolor{rowgrey} Valeur invalide / capteur HS & 0xFF (uint8), 0xFFFF (uint16), 0x7FFF (int16) & \coche{cvok3} & \champ{cvc3} \\
  Bits non utilisés & Forcés à 0 & \coche{cvok4} & \champ{cvc4} \\
  \bottomrule
\end{tabularx}

\subsection*{Exemples}
\begin{longtable}{L{3.2cm} C{1.7cm} C{1.4cm} L{4.2cm} L{2.0cm} L{3.6cm} L{5.4cm}}
  \toprule
  \textbf{Message} & \textbf{Octet(s)} & \textbf{Bits} & \textbf{Signal} & \textbf{Type} & \textbf{Codage (unité $\times$ facteur)} & \textbf{Plage, valeur invalide} \\
  \midrule
  \endhead
  \ex{WHEEL\_SPEED} & \ex{0 -- 1} & \ex{0 -- 15} & \ex{Vitesse roue} & \ex{uint16} & \ex{km/h $\times$ 100} & \ex{0 à 655,35 ; 0xFFFF = invalide} \\
  \ex{IMU\_DATA} & \ex{0 -- 1} & \ex{0 -- 15} & \ex{Accélération X} & \ex{int16} & \ex{g $\times$ 1000} & \ex{$\pm$32,7 g ; 0x7FFF = invalide} \\
  \ex{IMU\_DATA} & \ex{2 -- 3} & \ex{0 -- 15} & \ex{Accélération Y} & \ex{int16} & \ex{g $\times$ 1000} & \ex{$\pm$32,7 g ; 0x7FFF = invalide} \\
  \ex{IMU\_DATA} & \ex{4 -- 5} & \ex{0 -- 15} & \ex{Vitesse angulaire Z} & \ex{int16} & \ex{$^\circ$/s $\times$ 10} & \ex{$\pm$3276 $^\circ$/s} \\
  \ex{NODE\_STATUS} & \ex{0} & \ex{0} & \ex{Capteur vitesse OK} & \ex{bit} & \ex{1 = OK} & \ex{---} \\
  \ex{NODE\_STATUS} & \ex{0} & \ex{1} & \ex{Capteur température OK} & \ex{bit} & \ex{1 = OK} & \ex{---} \\
  \bottomrule
\end{longtable}

\noindent\footnotesize\textit{Comment on arrive au codage de WHEEL\_SPEED : plage 0 à 80 km/h, précision 0,01 km/h
$\rightarrow$ 8000 valeurs distinctes $\rightarrow$ 13 bits nécessaires $\rightarrow$ arrondi à 16 bits (2 octets), uint16 car toujours positif.}
\normalsize

\clearpage
\subsection*{À remplir}
\begin{longtable}{L{3.2cm} C{1.7cm} C{1.4cm} L{4.2cm} L{2.0cm} L{3.6cm} L{5.4cm}}
  \toprule
  \textbf{Message} & \textbf{Octet(s)} & \textbf{Bits} & \textbf{Signal} & \textbf{Type} & \textbf{Codage (unité $\times$ facteur)} & \textbf{Plage, valeur invalide} \\
  \midrule
  \endhead

  \champ{k1msg} & \champ{k1oct} & \champ{k1bit} & \champ{k1sig} & \champ{k1typ} & \champ{k1cod} & \champ{k1pl} \\
  \champ{k2msg} & \champ{k2oct} & \champ{k2bit} & \champ{k2sig} & \champ{k2typ} & \champ{k2cod} & \champ{k2pl} \\
  \champ{k3msg} & \champ{k3oct} & \champ{k3bit} & \champ{k3sig} & \champ{k3typ} & \champ{k3cod} & \champ{k3pl} \\
  \champ{k4msg} & \champ{k4oct} & \champ{k4bit} & \champ{k4sig} & \champ{k4typ} & \champ{k4cod} & \champ{k4pl} \\
  \champ{k5msg} & \champ{k5oct} & \champ{k5bit} & \champ{k5sig} & \champ{k5typ} & \champ{k5cod} & \champ{k5pl} \\
  \champ{k6msg} & \champ{k6oct} & \champ{k6bit} & \champ{k6sig} & \champ{k6typ} & \champ{k6cod} & \champ{k6pl} \\
  \champ{k7msg} & \champ{k7oct} & \champ{k7bit} & \champ{k7sig} & \champ{k7typ} & \champ{k7cod} & \champ{k7pl} \\
  \champ{k8msg} & \champ{k8oct} & \champ{k8bit} & \champ{k8sig} & \champ{k8typ} & \champ{k8cod} & \champ{k8pl} \\
  \champ{k9msg} & \champ{k9oct} & \champ{k9bit} & \champ{k9sig} & \champ{k9typ} & \champ{k9cod} & \champ{k9pl} \\
  \champ{k10msg} & \champ{k10oct} & \champ{k10bit} & \champ{k10sig} & \champ{k10typ} & \champ{k10cod} & \champ{k10pl} \\
  \champ{k11msg} & \champ{k11oct} & \champ{k11bit} & \champ{k11sig} & \champ{k11typ} & \champ{k11cod} & \champ{k11pl} \\
  \champ{k12msg} & \champ{k12oct} & \champ{k12bit} & \champ{k12sig} & \champ{k12typ} & \champ{k12cod} & \champ{k12pl} \\
  \champ{k13msg} & \champ{k13oct} & \champ{k13bit} & \champ{k13sig} & \champ{k13typ} & \champ{k13cod} & \champ{k13pl} \\
  \champ{k14msg} & \champ{k14oct} & \champ{k14bit} & \champ{k14sig} & \champ{k14typ} & \champ{k14cod} & \champ{k14pl} \\
  \champ{k15msg} & \champ{k15oct} & \champ{k15bit} & \champ{k15sig} & \champ{k15typ} & \champ{k15cod} & \champ{k15pl} \\
  \champ{k16msg} & \champ{k16oct} & \champ{k16bit} & \champ{k16sig} & \champ{k16typ} & \champ{k16cod} & \champ{k16pl} \\
  \champ{k17msg} & \champ{k17oct} & \champ{k17bit} & \champ{k17sig} & \champ{k17typ} & \champ{k17cod} & \champ{k17pl} \\
  \champ{k18msg} & \champ{k18oct} & \champ{k18bit} & \champ{k18sig} & \champ{k18typ} & \champ{k18cod} & \champ{k18pl} \\
  \champ{k19msg} & \champ{k19oct} & \champ{k19bit} & \champ{k19sig} & \champ{k19typ} & \champ{k19cod} & \champ{k19pl} \\
  \champ{k20msg} & \champ{k20oct} & \champ{k20bit} & \champ{k20sig} & \champ{k20typ} & \champ{k20cod} & \champ{k20pl} \\
  \champ{k21msg} & \champ{k21oct} & \champ{k21bit} & \champ{k21sig} & \champ{k21typ} & \champ{k21cod} & \champ{k21pl} \\
  \champ{k22msg} & \champ{k22oct} & \champ{k22bit} & \champ{k22sig} & \champ{k22typ} & \champ{k22cod} & \champ{k22pl} \\
  \champ{k23msg} & \champ{k23oct} & \champ{k23bit} & \champ{k23sig} & \champ{k23typ} & \champ{k23cod} & \champ{k23pl} \\
  \champ{k24msg} & \champ{k24oct} & \champ{k24bit} & \champ{k24sig} & \champ{k24typ} & \champ{k24cod} & \champ{k24pl} \\
  \champ{k25msg} & \champ{k25oct} & \champ{k25bit} & \champ{k25sig} & \champ{k25typ} & \champ{k25cod} & \champ{k25pl} \\
  \champ{k26msg} & \champ{k26oct} & \champ{k26bit} & \champ{k26sig} & \champ{k26typ} & \champ{k26cod} & \champ{k26pl} \\
  \champ{k27msg} & \champ{k27oct} & \champ{k27bit} & \champ{k27sig} & \champ{k27typ} & \champ{k27cod} & \champ{k27pl} \\
  \champ{k28msg} & \champ{k28oct} & \champ{k28bit} & \champ{k28sig} & \champ{k28typ} & \champ{k28cod} & \champ{k28pl} \\
  \champ{k29msg} & \champ{k29oct} & \champ{k29bit} & \champ{k29sig} & \champ{k29typ} & \champ{k29cod} & \champ{k29pl} \\
  \champ{k30msg} & \champ{k30oct} & \champ{k30bit} & \champ{k30sig} & \champ{k30typ} & \champ{k30cod} & \champ{k30pl} \\
  \champ{k31msg} & \champ{k31oct} & \champ{k31bit} & \champ{k31sig} & \champ{k31typ} & \champ{k31cod} & \champ{k31pl} \\
  \champ{k32msg} & \champ{k32oct} & \champ{k32bit} & \champ{k32sig} & \champ{k32typ} & \champ{k32cod} & \champ{k32pl} \\
  \champ{k33msg} & \champ{k33oct} & \champ{k33bit} & \champ{k33sig} & \champ{k33typ} & \champ{k33cod} & \champ{k33pl} \\
  \champ{k34msg} & \champ{k34oct} & \champ{k34bit} & \champ{k34sig} & \champ{k34typ} & \champ{k34cod} & \champ{k34pl} \\
  \champ{k35msg} & \champ{k35oct} & \champ{k35bit} & \champ{k35sig} & \champ{k35typ} & \champ{k35cod} & \champ{k35pl} \\
  \champ{k36msg} & \champ{k36oct} & \champ{k36bit} & \champ{k36sig} & \champ{k36typ} & \champ{k36cod} & \champ{k36pl} \\
  \champ{k37msg} & \champ{k37oct} & \champ{k37bit} & \champ{k37sig} & \champ{k37typ} & \champ{k37cod} & \champ{k37pl} \\
  \champ{k38msg} & \champ{k38oct} & \champ{k38bit} & \champ{k38sig} & \champ{k38typ} & \champ{k38cod} & \champ{k38pl} \\
  \champ{k39msg} & \champ{k39oct} & \champ{k39bit} & \champ{k39sig} & \champ{k39typ} & \champ{k39cod} & \champ{k39pl} \\
  \champ{k40msg} & \champ{k40oct} & \champ{k40bit} & \champ{k40sig} & \champ{k40typ} & \champ{k40cod} & \champ{k40pl} \\
  \champ{k41msg} & \champ{k41oct} & \champ{k41bit} & \champ{k41sig} & \champ{k41typ} & \champ{k41cod} & \champ{k41pl} \\
  \champ{k42msg} & \champ{k42oct} & \champ{k42bit} & \champ{k42sig} & \champ{k42typ} & \champ{k42cod} & \champ{k42pl} \\
  \champ{k43msg} & \champ{k43oct} & \champ{k43bit} & \champ{k43sig} & \champ{k43typ} & \champ{k43cod} & \champ{k43pl} \\
  \champ{k44msg} & \champ{k44oct} & \champ{k44bit} & \champ{k44sig} & \champ{k44typ} & \champ{k44cod} & \champ{k44pl} \\
  \champ{k45msg} & \champ{k45oct} & \champ{k45bit} & \champ{k45sig} & \champ{k45typ} & \champ{k45cod} & \champ{k45pl} \\
  \champ{k46msg} & \champ{k46oct} & \champ{k46bit} & \champ{k46sig} & \champ{k46typ} & \champ{k46cod} & \champ{k46pl} \\
  \champ{k47msg} & \champ{k47oct} & \champ{k47bit} & \champ{k47sig} & \champ{k47typ} & \champ{k47cod} & \champ{k47pl} \\
  \champ{k48msg} & \champ{k48oct} & \champ{k48bit} & \champ{k48sig} & \champ{k48typ} & \champ{k48cod} & \champ{k48pl} \\
  \champ{k49msg} & \champ{k49oct} & \champ{k49bit} & \champ{k49sig} & \champ{k49typ} & \champ{k49cod} & \champ{k49pl} \\
  \champ{k50msg} & \champ{k50oct} & \champ{k50bit} & \champ{k50sig} & \champ{k50typ} & \champ{k50cod} & \champ{k50pl} \\
  \champ{k51msg} & \champ{k51oct} & \champ{k51bit} & \champ{k51sig} & \champ{k51typ} & \champ{k51cod} & \champ{k51pl} \\
  \champ{k52msg} & \champ{k52oct} & \champ{k52bit} & \champ{k52sig} & \champ{k52typ} & \champ{k52cod} & \champ{k52pl} \\
  \champ{k53msg} & \champ{k53oct} & \champ{k53bit} & \champ{k53sig} & \champ{k53typ} & \champ{k53cod} & \champ{k53pl} \\
  \champ{k54msg} & \champ{k54oct} & \champ{k54bit} & \champ{k54sig} & \champ{k54typ} & \champ{k54cod} & \champ{k54pl} \\
  \champ{k55msg} & \champ{k55oct} & \champ{k55bit} & \champ{k55sig} & \champ{k55typ} & \champ{k55cod} & \champ{k55pl} \\
  \champ{k56msg} & \champ{k56oct} & \champ{k56bit} & \champ{k56sig} & \champ{k56typ} & \champ{k56cod} & \champ{k56pl} \\
  \bottomrule
\end{longtable}
\setlength{\tabcolsep}{6pt}
\portrait


\section{Contraintes électriques et sécurité}

Si vous ne savez pas encore, écrivez « à voir » — c'est une réponse utile.

\begin{tabularx}{\textwidth}{L{5.6cm} X}
  \toprule
  \textbf{Question} & \textbf{Votre réponse} \\
  \midrule

  Tension du réseau de bord\newline{\footnotesize\textit{5 V / 12 V / 24 V / autre}} & \champML{cn1} \\
  Un 5 V est-il disponible sur vos cartes ?\newline{\footnotesize\textit{conditionne le choix du transceiver (voir section 3)}} & \champML{cn2} \\
  Électronique de puissance présente ?\newline{\footnotesize\textit{onduleur, variateur, gros moteur}} & \champML{cn3} \\
  Sources de bruit CEM identifiées\newline{\footnotesize\textit{commutation, alimentation à découpage}} & \champML{cn4} \\
  Isolation galvanique nécessaire ?\newline{\footnotesize\textit{entre quels sous-systèmes ?}} & \champML{cn5} \\
  Système d'arrêt d'urgence\newline{\footnotesize\textit{existe-t-il ? passe-t-il par le bus ?}} & \champML{cn6} \\
  Comportement attendu si le bus tombe\newline{\footnotesize\textit{mode dégradé, arrêt sécurisé, valeurs par défaut}} & \champML{cn7} \\
  Timeout de perte d'un message critique\newline{\footnotesize\textit{ex. absence pendant 500 ms = arrêt}} & \champML{cn8} \\
  Longueur approximative du bus\newline{\footnotesize\textit{en mètres, d'un bout à l'autre}} & \champML{cn9} \\
  Connecteur souhaité ou imposé\newline{\footnotesize\textit{si vous avez déjà une contrainte}} & \champML{cn10} \\
  Contraintes mécaniques\newline{\footnotesize\textit{place, vibrations, étanchéité}} & \champML{cn11} \\
  Précautions de routage / blindage & \champML{cn12} \\
  \bottomrule
\end{tabularx}


\clearpage
\section{Points ouverts}

Questions sans réponse à la date de cette version. C'est normal d'en avoir : mieux vaut un point
ouvert écrit qu'une hypothèse silencieuse.

\subsection*{Exemple}
\begin{tabularx}{\textwidth}{C{0.8cm} X L{2.8cm} L{2.0cm} L{1.8cm}}
  \toprule
  \textbf{N\textsuperscript{o}} & \textbf{Question} & \textbf{Qui répond} & \textbf{Échéance} & \textbf{Statut} \\
  \midrule
  \ex{1} & \ex{Le flux des caméras passe-t-il par le bus ou par une liaison dédiée ?} & \ex{Équipe vision} & \ex{Séance 3} & \ex{Ouvert} \\
  \bottomrule
\end{tabularx}

\subsection*{À remplir}
\begin{tabularx}{\textwidth}{C{0.8cm} X L{2.8cm} L{2.0cm} L{1.8cm}}
  \toprule
  \textbf{N\textsuperscript{o}} & \textbf{Question} & \textbf{Qui répond} & \textbf{Échéance} & \textbf{Statut} \\
  \midrule

  1 & \champML{op1q} & \champML{op1w} & \champML{op1e} & \champML{op1s} \\
  2 & \champML{op2q} & \champML{op2w} & \champML{op2e} & \champML{op2s} \\
  3 & \champML{op3q} & \champML{op3w} & \champML{op3e} & \champML{op3s} \\
  4 & \champML{op4q} & \champML{op4w} & \champML{op4e} & \champML{op4s} \\
  5 & \champML{op5q} & \champML{op5w} & \champML{op5e} & \champML{op5s} \\
  6 & \champML{op6q} & \champML{op6w} & \champML{op6e} & \champML{op6s} \\
  7 & \champML{op7q} & \champML{op7w} & \champML{op7e} & \champML{op7s} \\
  8 & \champML{op8q} & \champML{op8w} & \champML{op8e} & \champML{op8s} \\
  9 & \champML{op9q} & \champML{op9w} & \champML{op9e} & \champML{op9s} \\
  10 & \champML{op10q} & \champML{op10w} & \champML{op10e} & \champML{op10s} \\
  11 & \champML{op11q} & \champML{op11w} & \champML{op11e} & \champML{op11s} \\
  12 & \champML{op12q} & \champML{op12w} & \champML{op12e} & \champML{op12s} \\
  \bottomrule
\end{tabularx}

\vfill
\noindent\footnotesize\textit{Merci d'avoir rempli ce document, même partiellement.
Enregistrez le PDF et renvoyez-le à hugo.mathoux@ensea.fr.}
\normalsize

\end{Form}
\end{document}